% Bagean: sajarah
% Bab: biografi
% Pérangan: kurt-goedel

\documentclass[../../../include/open-logic-section]{subfiles}

\begin{document}

\olfileid{his}{bio}{god}

\olsection{Kurt G\"odel}

\olphoto{goedel-kurt}{Kurt G{\"o}del}

Kurt G{\"o}del (\textsc{ger}-dle) lair tanggal 28 April
1906 ing Br{\"u}nn, Kakaisaran Austria-Hongaria
(saiki Brno ing Républik Ceko). Amarga wataké
seneng takon lan pinter, Kurtele cilik kerep
diarani ``Der kleine Herr Warum'' (Pak Cilik
Yagéné) déning kulawargané. Dhèwèké unggul ing
pasinaon wiwit sekolah dhasar; mung ing
matématika bijiné tau ora paling dhuwur.
G{\"o}del kerep ora mlebu sekolah amarga
kesehatané kurang apik lan ora diwajibaké
melu olah raga. Nalika isih cilik, dhèwèké
didiagnosis ngalami demam rematik. Sajeroning
uripé, dhèwèké percaya yèn penyakit iku
nyebabaké karusakan permanèn ing jantungé,
senajan pamriksaan mèdhis nyatakaké kosok
baliné.

G{\"o}del wiwit kuliah ing Universitas Wina
ing taun 1924 lan ngrampungaké studi
doktoral ing~1929. Wiwitané dhèwèké arep
nyinaoni fisika, nanging minaté enggal
ngalih menyang matématika, mligi logika,
antara liya amarga pangaruh filsuf
Rudolf Carnap. Disertasiné, kang ditulis
ing sangisoré bimbingan Hans Hahn,
mbuktèkaké teorema kelengkapan logika
predikat orde kapisan kanthi identitas
\citep{Godel1929}. Mung setaun sawisé iku,
dhèwèké olèh asilé kang paling misuwur---
teorema ketaklengkapan kapisan lan kapindho
(kababar ing \citealt{Godel1931}).
Nalika ing Wina, G{\"o}del aktif banget
ing Lingkaran Wina, sawijining kelompok
filsuf kang mikir kanthi cara ilmiah lan
kalebu Carnap; karya Carnap utamané
kapengaruh déning asil-asilé G{\"o}del.

Ing taun 1938, G\"odel omah-omah karo Adele
Nimbursky. Wong tuwané ora seneng:
Adele ora mung luwih tuwa nem taun lan
wis pegatan, nanging uga nyambut gawé
minangka panari ing klub wengi.
Nanging tekanan sosial ora mengaruhi
G{\"o}del, lan kekaroné tetep urip
bebrayan kanthi bungah nganti dhèwèké séda.

Sawisé Jerman Nazi nggabungaké Austria
ing taun 1938, G{\"o}del lan Adele
pindhah menyang Amérika Sarékat, lan
dhèwèké nampa jabatan ing Institute
for Advanced Study ing Princeton,
New Jersey. Senajan rada meneng lan nduwèni
watak kang ora lumrah, wektu G{\"o}del
ing Princeton kebak kerja bareng lan
ngasilaké akèh karya. Dhèwèké
nerbitaké ésai ngenani téori himpunan,
filsafat, lan fisika. Mliginé,
dhèwèké dadi kanca cedhak banget
karo Albert Einstein, kolegané ing IAS.

Ing taun-taun pungkasan uripé,
kesehatan mental G{\"o}del saya
mudhun. Amarga garwané dirawat
ing rumah sakit ing taun 1977,
garwané ora bisa masakaké panganan
kanggo dhèwèké. Sawisé ngalami
masalah kesehatan mental sajroning
uripé, dhèwèké banjur kena
paranoia. Amarga wedi banget
diracuni, G{\"o}del ora gelem
mangan. Dhèwèké séda amarga
kaliren tanggal 14 Januari 1978
ing Princeton.

\begin{reading}
Kanggo biografi jangkep ngenani
uripé G{\"o}del, delengen
\citet{Dawson1997}. Kanggo tulisan
biografi liya lan ésai ngenani
sumbangané G{\"o}del marang
logika lan filsafat, delengen
\citet{Wang1990}, \citet{Baaz2011},
\citet{Takeuti2003}, lan
\citet{Sigmund2007}.

Tèsis doktor G{\"o}del kasedhiya
ing basa Jerman asliné
\citep{Godel1929}. Tèks asli
teorema-teorema ketaklengkapan
ana ing \citep{Godel1931}.
Kabèh tulisan G\"odel kang
wis utawa durung kababar,
uga pilihan surat-menyuraté,
kasedhiya ing basa Inggris
ing \emph{Collected Papers}
\citet{Godel1986,Godel1990}.

Kanggo pamaparan rinci ngenani
teorema-teorema ketaklengkapan
G{\"o}del, delengen
\citet{Smith2013}. Kanggo
rembugan filosofis kang ora
formal ngenani teorema-teorema
G{\"o}del, delengen podcast
Mark Linsenmayer
\citep{Linsenmayer2014}.
\end{reading}

\end{document}
